\chapter{Literature Review and State-of-the-Art}
\label{ch:sota}

Environmental perception in autonomous driving has evolved substantially since
the early rule-based systems, progressing from monocular detection pipelines on
individual vehicles toward fully learned, multi-agent systems capable of
reasoning about shared perceptual information in real time. Modern cooperative
perception sits at the intersection of three distinct research fields that have
individually achieved significant maturity: stereo-based 3D perception, deep
learning for object detection, and collaborative perception networks. This
chapter surveys the current state of the art across each of these three domains,
tracing the trajectory of each field and identifying where the boundaries of
current methods lie.

\section{Stereo-Based 3D Perception}
\label{sec:stereo3d}

Autonomous driving platforms require a real-time, metric three-dimensional
representation of their surroundings before any higher-level task such as object
tracking, trajectory prediction, or path planning can proceed reliably. How that
representation is obtained from onboard sensors is a non-trivial design choice,
and one that has shifted considerably as the field has matured. This section
examines the principal sensing paradigms in autonomous perception, then traces
the algorithmic evolution of stereo depth estimation methods.

\subsection{Sensor Modalities}

To capture the spatial layout of an environment, modern autonomous perception
systems primarily utilise one of three core sensor configurations: active LiDAR arrays, passive monocular cameras, or passive
stereo cameras. Each modality presents a fundamentally distinct engineering
compromise across financial cost, hardware reliability, semantic density, and
geometric depth accuracy~\cite{valverde2025}.

Active sensing via LiDAR has long served as the benchmark for metric 3D
perception. By emitting localised infrared laser pulses and calculating the
direct Time-of-Flight (ToF) of their reflections, LiDAR sensors generate sparse
but highly accurate point clouds that are largely invariant to ambient lighting
conditions or shadows. This direct coordinate measurement eliminates downstream
geometric ambiguity, providing high metric precision at ranges exceeding one
hundred metres. Despite these geometric advantages, active LiDAR sensing suffers
from severe operational constraints: high-resolution automotive units remain
prohibitively expensive for mass-market deployment, carry significant power
consumption, rely on mechanical or solid-state components vulnerable to physical
wear, and produce point clouds inherently devoid of semantic texture, rendering
them incapable of capturing colour, reading signage, or interpreting traffic
state.

Conversely, monocular camera systems offer the most affordable and
computationally streamlined alternative for automotive perception. Single cameras
provide dense, high-resolution RGB images that contain rich semantic textures,
allowing deep neural networks to classify objects, detect lane boundaries, and
interpret complex scene contexts. However, extracting 3D metric coordinates from
a single 2D image plane is a mathematically ill-posed inverse problem. Because
the projection from a 3D world space to a 2D pixel grid discards the depth
dimension entirely, monocular systems suffer from fundamental scale
ambiguity~\cite{huang2025monocular}. To estimate distance, monocular deep
learning models must rely heavily on statistical priors and contextual cues
learned during training, such as assuming a vehicle's standard height or
ground-plane consistency. Consequently, when confronted with out-of-distribution
obstacles, unusual vehicle profiles, or severe ego-vehicle pitch, these models
degrade substantially, with published depth errors typically falling in the
15--30\% range and accuracy degrading further beyond 50 metres as perspective
cues become insufficient to resolve absolute scale~\cite{huang2025}.

Stereo vision recovers absolute metric depth through passive geometric
triangulation across two synchronised, horizontally offset cameras separated by a
fixed known baseline, requiring no learned priors about object appearance or
scale. This geometry preserves the dense semantic richness and low hardware cost
of camera-based sensing while establishing a physically grounded foundation for
depth recovery. Unlike learned monocular approaches, stereo triangulation is
distribution-agnostic: it applies equally to any rigid object in the scene
regardless of whether that object class appeared in the depth model's training
data, making it a practical and scalable baseline for single-agent 3D perception
in cost-constrained deployments~\cite{huang2025}.

\subsection{Stereo Depth Estimation}
\label{sec:stereodepth}

Extracting reliable depth from a stereo pair requires solving a precise geometric
problem that has driven decades of algorithmic development. The foundational
objective of a computational stereo engine is to solve the correspondence
problem: identifying the exact spatial location of a point in a secondary camera
frame that matches a given pixel in the primary camera frame. By assuming a
rectified camera arrangement where the optical axes are parallel, the search
space for a matching pixel can be mathematically constrained from a
two-dimensional grid to a one-dimensional horizontal epipolar line. The
horizontal pixel displacement between these matched points represents the
disparity, which inversely determines the absolute metric depth of the scene
relative to the camera's baseline and focal length.

Historically, the computer vision literature divided dense disparity estimation
into local and global optimisation paradigms~\cite{scharstein2002}. Local
methods, such as Block Matching (BM), compute correspondence costs over isolated
support windows using pixel-level intensity metrics. While computationally
efficient, local matching degrades significantly in textureless regions or under
uniform illumination. Global matching paradigms address these spatial ambiguities
by introducing mathematical energy minimisation functions. Algorithms like
Semi-Global Matching (SGM)~\cite{hirschmuller2008} approximate global 2D
optimisation by aggregating dynamic programming paths across multiple 1D
orientations through the image grid. SGM balances per-pixel matching accuracy
with contextual smoothness constraints, making it a practical baseline for
real-time stereo matching in robotics and automotive applications. However, it
still exhibits severe optimisation failures when confronted with reflective
materials, fine structural boundaries, or repeated patterns.

The integration of deep learning drastically shifted the stereo landscape by
replacing hand-crafted cost functions with data-driven representations.
Influential deep architectures such as PSMNet~\cite{chang2018} introduced the
concept of the 3D visual cost volume. Instead of matching raw pixel intensities,
these networks employ Siamese Convolutional Neural Networks (CNNs) to extract
dense feature descriptors from both views, which are concatenated across a range
of predefined candidate disparities to construct a high-dimensional cost volume.
Volumetric 3D convolutions are then applied to aggregate spatial and disparity
context over the volume to regress a smooth disparity map.

To refine this volumetric approach, the field transitioned toward iterative
recurrent architectures, an approach exemplified by RAFT-Stereo~\cite{lipson2021}.
RAFT-Stereo extracts multi-scale feature maps, constructs a multi-scale 4D
correlation cost volume, and utilises a Gated Recurrent Unit (GRU) block to
iteratively look up and update the disparity field. While RAFT-Stereo achieved
unprecedented accuracy on fine structures across major benchmarks, maintaining
and fetching data from multi-scale 4D cost volumes still imposes significant
memory bottlenecks on embedded automotive hardware.

To completely bypass the computational overhead imposed by volumetric
representations, contemporary research has shifted toward cost-volume-free,
recurrent warping architectures, an approach exemplified by
WAFT-Stereo~\cite{wang2026}. Rather than building or maintaining a static,
memory-heavy cost structure, WAFT-Stereo frames stereo matching as a continuous
field transformation, iteratively warping feature maps from one view directly
onto the other and utilising localised field transforms to resolve geometric
misalignments progressively. By relying strictly on high-resolution warping
mechanisms rather than dense volumetric convolutions or correlation lookups, this
framework holds a top-ranking position on ETH3D, Middlebury, and KITTI across all
metrics, while running 1.8 to 6.7 times faster than competitive
methods~\cite{wang2026}, providing a highly scalable depth engine for downstream
detection tasks.

\Cref{tab:stereo} provides a comparative summary of this algorithmic evolution
across four representative methodologies, contrasting their core paradigms, error
metrics on the KITTI 2015 benchmark, operational latencies, and deployment
hardware configurations.

\begin{table}[ht]
\centering
\caption{Representative stereo depth estimation methods on the KITTI 2015 benchmark.}
\label{tab:stereo}
\renewcommand{\arraystretch}{1.6}
\small
\resizebox{\textwidth}{!}{%
\begin{tabular}{l l c c l}
\toprule
\thead{Method} & \thead{Paradigm} &
\thead{KITTI 2015 \\ D1-all} &
\thead{Latency} & \thead{Hardware} \\
\midrule
SGBM \cite{hirschmuller2008}
  & Classical path aggregation
  & $\sim$8--12\% & $<$25\,ms & Multi-core CPU \\[4pt]
PSMNet \cite{chang2018}
  & CNN + 3D cost volume
  & 2.32\% & $\sim$410\,ms & NVIDIA GTX 1080Ti \\[4pt]
RAFT-Stereo \cite{lipson2021}
  & Recurrent + 4D correlation volume
  & 2.42\% & $\sim$380\,ms & NVIDIA RTX 3090 \\[4pt]
WAFT-Stereo \cite{wang2026}
  & Recurrent + feature warping (no cost volume)
  & 1.28\% & $\sim$100\,ms & NVIDIA L40 \\
\bottomrule
\end{tabular}}
\end{table}





\section{Single-Vehicle 3D Object Detection}
\label{sec:single3d}

Once a dense depth map is generated from the ego-vehicle's stereo vision engine,
the perception pipeline must use this depth information to localise and classify
surrounding vehicles. Single-vehicle 3D object detection focuses on predicting a
bounding box defined by a seven-degree-of-freedom state vector: a 3D centre
coordinate $(x, y, z)$, physical dimensions (height, width, length), and an
orientation angle (yaw). Within the autonomous vehicle literature, deep
architectures handle this localisation task through two fundamentally distinct
structural strategies. The first is a multi-stage lifting paradigm, which relies
on an explicit 2D image-plane detector to isolate objects before projecting them
into 3D coordinates using calculated depth data. The second is an end-to-end
paradigm encompassing volumetric and transformer-based architectures, which
bypasses intermediate 2D detections entirely, directly mapping raw stereo inputs
to 3D spatial parameters. This section first evaluates the core 2D detection
architectures that serve as the front-end for lifting workflows, before examining
both multi-stage pseudo-LiDAR approaches and end-to-end direct regression
systems.

\subsection{2D Object Detection as a Front-End for 3D Lifting}

In multi-stage perception frameworks, a 2D image-plane detector serves as a
structural prerequisite: it must first isolate objects within the RGB image
plane, whose boundaries are then used to crop depth maps and define the spatial
frustum passed to the 3D lifting stage. Because any localisation error at this
stage directly bounds the accuracy of the entire pipeline, the choice of 2D
backbone carries significant architectural weight. The literature primarily
offers three configurations for image-plane localisation: classical two-stage
proposal networks, modern single-shot convolutional networks, and real-time
end-to-end vision transformers.

The canonical two-stage detection baseline is Faster R-CNN~\cite{ren2015}. This
architecture established the foundational paradigm of separating localisation into
a sequential processing pipeline. The first stage employs a dedicated Region
Proposal Network (RPN) to scan dense feature maps and generate coarse bounding
box candidates. The second stage then utilises region-of-interest (RoI) pooling
layers to extract these features, passing them to downstream heads for final
categorical classification and bounding box regression. While Faster R-CNN
consistently achieved superior localisation precision and robust handling of
multi-scale objects compared to single-stage methods of the same period, its
multi-stage design introduces computational bottlenecks on embedded hardware,
making real-time deployment on automotive hardware challenging.

To address these latency constraints, the field shifted toward one-stage,
single-shot detectors, exemplified by the YOLO family (You Only Look
Once)~\cite{redmon2016}. One-stage architectures eliminate the region proposal
step entirely, framing object localisation as a direct, end-to-end regression
problem over a dense spatial grid. While early iterations relied on rigid,
hand-crafted anchor boxes that struggled with objects of varying scales and
aspect ratios, the family has evolved continuously, progressively adopting
anchor-free regression heads and increasingly expressive backbone designs. The
most recent iteration, YOLOv12, builds on this trajectory by introducing a
lightweight area attention module and a Residual Efficient Layer Aggregation
Network (R-ELAN), achieving incremental mAP gains over its predecessor while
matching its inference speed on standard GPU hardware~\cite{tian2025}, continuing
to narrow the historical trade-off between inference speed and detection accuracy.

Beyond convolutional architectures, contemporary multi-stage frameworks have
increasingly adopted Vision Transformers (ViTs), specifically the DETR (DEtection
TRansformer) lineage, as their localisation engine~\cite{carion2020}. While
convolutional models rely on localised receptive fields and hand-crafted
post-processing heuristics like Non-Maximum Suppression (NMS) to clear out
redundant, overlapping boxes, DETR models frame 2D localisation as a direct
set-prediction problem. Utilising multi-head self-attention mechanisms, these
networks capture global contextual dependencies and long-range cross-object
interactions across the entire image plane, matching learned object queries to
ground-truth boxes via bipartite matching. With modern adaptations like
RT-DETR~\cite{lv2024} drastically optimising the computational latency of these
attention maps, transformer-based frameworks produce exceptionally clean,
precisely bounded region crops without spatial suppression artefacts, a property
that directly improves the quality of the depth frustum passed to the 3D lifting
stage.

\Cref{tab:detectors} compares the three architectural families across detection
accuracy, inference latency, and primary deployment context. Latency values for
YOLOv12 and RT-DETR are measured using TensorRT FP16 (TRT FP16) optimisation on
an NVIDIA T4 GPU; Faster R-CNN latency is measured in standard PyTorch on a server
GPU.

\begin{table}[ht]
\centering
\caption{Comparison of major 2D object detection models on the COCO benchmark.}
\label{tab:detectors}
\renewcommand{\arraystretch}{1.6}
\small
\resizebox{\textwidth}{!}{%
\begin{tabular}{l l c c l}
\toprule
\thead{Model} & \thead{Family} &
\thead{COCO \\ mAP@0.5:0.95} &
\thead{Latency} & \thead{Design Priority} \\
\midrule
Faster R-CNN \cite{ren2015}
  & Two-stage proposal
  & $\sim$37\% & 100--200\,ms & Precision-first, latency-tolerant \\[4pt]
YOLOv12 \cite{tian2025}
  & One-stage anchor-free + attention
  & $\sim$51.8\% & $\sim$2.5\,ms TRT FP16 & Speed-first, edge deployment \\[4pt]
RT-DETR \cite{lv2024}
  & Transformer set prediction
  & $\sim$53.0\% & $\sim$9.3\,ms TRT FP16 & Transformer precision at real-time speed \\
\bottomrule
\end{tabular}}
\end{table}

\subsection{Stereo-Based 3D Object Detection}

While image-plane analysis extracts rich two-dimensional semantic maps,
autonomous vehicles require structural reasoning in native three-dimensional
coordinate frames. Resolving the extra spatial dimension from stereo camera
setups has produced two structurally distinct paradigms in the literature, driven
by a fundamental tension between geometric interpretability and end-to-end
computational efficiency.

\paragraph{a. Two-Stage Lifting and Pseudo-LiDAR Approaches.}
The two-stage lifting paradigm connects image-plane processing to 3D geometric
interpretation by decoupling the pipeline into explicit depth generation and
point-cloud detection components. This approach was popularised through the
Pseudo-LiDAR framework~\cite{wang2019pseudolidar}, which applies back-projection
to estimated disparity maps. Utilising camera calibration matrices, the network
converts every pixel with a known depth value directly into a concrete 3D point
cloud within the physical world coordinate frame. Once this coordinate
transformation occurs, the generated Pseudo-LiDAR point cloud mimics the
structural format of LiDAR point clouds. Consequently, these architectures can
directly deploy existing 3D point-cloud detection models, such as
PointRCNN~\cite{shi2019}, to group spatial points and regress 3D bounding boxes.
Subsequent iterations like Pseudo-LiDAR++ refined this by using sparse physical
measurements to correct depth errors, which naturally scale quadratically with
actual physical distance~\cite{you2020}. While this explicit separation makes the
overall system highly interpretable and modular, it introduces substantial
deployment overhead because the latency of the stereo depth estimation engine and
the point-cloud processing backend stack sequentially.

\paragraph{b. End-to-End Volumetric and Transformer Architectures.}
To eliminate the computational bottlenecks and error propagation of multi-stage
pipelines, contemporary research focuses on end-to-end architectures that bypass
explicit point-cloud reconstruction. Early approaches within this paradigm built
explicit geometric volumes directly from stereo features, training unified
networks to jointly estimate depth and detect objects without intermediate
pseudo-LiDAR conversion. Early geometry-aware methods such as
DSGN~\cite{chen2020dsgn} constructed plane-sweep stereo volumes and processed them
with 3D convolutions to jointly reason about depth and object presence within a
unified feature space. Building on this geometric grounding,
LIGA-Stereo~\cite{guo2021} introduced LiDAR-guided supervision to strengthen
stereo feature learning, achieving high localisation accuracy on the KITTI
benchmark. However, the dense geometric volumes required by these architectures
impose significant memory and latency overhead, making real-time deployment on
automotive hardware difficult.

To overcome this latency barrier, contemporary frameworks like
StereoDETR~\cite{mu2025} decouple the network into two distinct branches: a fast
monocular DETR branch that regresses object scale, orientation, and spatial
queries, and a highly efficient stereo branch that leverages low-cost,
multi-scale disparity features to predict object-level depth maps. Coupled through
a differentiable depth sampling strategy, and handling occlusion through a
constrained supervision strategy that requires no extra annotations, StereoDETR
achieves real-time inference and is the first stereo-based method to surpass
monocular approaches in speed, while establishing competitive state-of-the-art
accuracy on the KITTI benchmark.

\section{Collaborative Perception via V2V Communication}
\label{sec:collab}

While modern single-vehicle 3D perception architectures have achieved strong
localisation accuracy, they remain inherently constrained by the physical limits
of a single perspective. Ego-vehicle sensors are fundamentally vulnerable to
long-range geometric degradation and critical blind spots caused by line-of-sight
occlusions.

Vehicle-to-Vehicle (V2V) cooperative perception addresses these single-agent
limitations by combining spatial data from multiple connected vehicles. By
sharing perception data across different vantage points, vehicles can mutually
resolve occlusions and extend their effective detection ranges. The primary
architectural variation among cooperative networks lies in the abstraction level
of the shared data, which is universally categorised into early, intermediate,
and late fusion paradigms. This section first examines these three fusion
paradigms, then traces the algorithmic milestones through which the field has
evolved.

\subsection{Fusion Paradigms}

The point within a neural network pipeline where data from neighbouring vehicles
is aggregated dictates the system's accuracy, bandwidth efficiency, and practical
deployment feasibility. The literature organises these methods into three distinct
approaches based on the level of information abstraction.

\paragraph{Early Fusion (Raw Data Aggregation).}
In early fusion pipelines, vehicles transmit raw, unprocessed sensor outputs, such
as stereo-derived dense depth maps or uncompressed point clouds, directly to the
ego-vehicle. The receiving agent aligns these raw inputs into its own spatial
coordinate system and processes the combined dense representation through a single
3D detection backbone. While early fusion preserves complete spatial features and
yields high theoretical localisation accuracy, it is highly impractical for
real-world vehicular setups. Transmitting raw, high-resolution visual data
generates massive bandwidth demands that easily saturate direct communication
channels, leading to severe transmission latency and data loss.

\paragraph{Intermediate Fusion (Feature-Level Aggregation).}
To balance bandwidth and performance, intermediate fusion strategies extract
multi-scale neural network feature maps, such as Bird's-Eye-View tensors, on
individual vehicles before transmission. Neighbouring vehicles exchange these
compressed latent features, which are then fused using spatial attention
mechanisms or graph neural networks. While intermediate fusion reduces bandwidth
relative to raw data, it introduces architectural coupling: cooperating vehicles
must share compatible backbone designs to ensure exchanged feature tensors align
semantically and can be combined by the downstream detection heads.

\paragraph{Late Fusion (Object-Level Aggregation).}
Late fusion operates at the highest level of data abstraction, completely
decoupling the individual perception pipelines of cooperating vehicles. Each
vehicle runs its own independent 3D object detection engine to generate a
finalised list of 3D bounding boxes, complete with spatial coordinates,
dimensions, orientations, and class confidence scores. Vehicles then transmit only
these low-dimensional object-level parameters across the V2V link. The receiving
ego-vehicle applies spatial transformation matrices to project the incoming
bounding boxes into its own coordinate frame, subsequently deploying clustering
algorithms, such as Intersection-over-Union matching or Kalman filtering, to fuse
duplicate detections and fill occluded blind spots. Because it transfers only
lightweight bounding box coordinates, late fusion minimises communication
bandwidth requirements and provides architectural flexibility, allowing
heterogeneous vehicle types with entirely different sensor setups and internal
detection backbones to collaborate without requiring shared neural network layers.

\subsection{Evolution of Collaborative Perception Systems}

To understand the current state of the art in multi-vehicle cooperative
perception, it is necessary to trace the algorithmic milestones that shaped the
field. The literature has evolved through a continuous cycle of establishing
fusion paradigms, identifying communication bottlenecks, and addressing the gap
between synthetic simulation and real-world deployment.

The foundational shift toward deep collaborative perception began with
V2VNet~\cite{wang2020v2vnet}. Prior to this work, cooperative approaches relied on
raw data sharing or disconnected object-level exchanges. V2VNet established the
core intermediate fusion paradigm, proving that vehicles could compress spatial
information into latent neural network representations and aggregate them using
graph neural networks. While it achieved significant accuracy gains over
single-vehicle baselines, V2VNet was evaluated purely on private synthetic
datasets where communication was assumed to be flawless and unconstrained.

This reliance on idealised conditions exposed a major structural void: the lack of
standardised benchmarking. OPV2V resolved this by introducing the first
open-source multi-vehicle cooperative perception benchmark and simulation
dataset~\cite{xu2022opv2v}. By providing a unified evaluation framework, OPV2V
allowed researchers to directly compare different fusion strategies under
identical conditions. Crucially, this benchmark exposed that dense intermediate
feature sharing incurred an unsustainably high communication bandwidth cost,
identifying data volume as the primary bottleneck at the time.

To alleviate this communication burden, Where2comm introduced the concept of
spatial communication efficiency~\cite{hu2022}. Recognising that large portions of
a vehicle's feature map contain redundant or empty spatial information, Where2comm
utilised a spatial confidence map to assess the perceptual informativeness of each
spatial region. Vehicles used this mechanism to dynamically select and transmit
only the most critical, high-entropy perception features, demonstrating that
communication volume could be drastically reduced without degrading downstream 3D
detection accuracy. CoBEVT generalised and scaled this selection mechanism by
introducing a unified transformer architecture~\cite{xu2022cobevt}. Instead of
relying on confidence-based spatial selection alone, CoBEVT designed a multi-agent
Bird's-Eye-View transformer that leveraged sparse axial self-attention, allowing
the ego-vehicle to perform context-aware spatial fusion across multiple connected
agents simultaneously while maintaining an efficient computational profile.

Despite these rapid architectural advancements, nearly all cooperative perception
research remained confined to synthetic simulators. V2V4Real confronted this
systemic limitation by releasing the first large-scale real-world cooperative
dataset collected by a physical connected vehicle fleet~\cite{xu2023v2v4real}.
V2V4Real exposed and quantified a significant sim-to-real gap, revealing that
real-world localisation errors, temporal asynchrony, and hardware heterogeneity
caused severe spatial misalignment that drastically degraded the performance of
tightly coupled intermediate feature fusion networks. Most recently, the
TUMTraf-V2X dataset extended the boundary of deployment-grade
evaluation~\cite{zimmer2024}. Collecting 2,000 labelled point clouds and 5,000
labelled images from five roadside and four onboard sensors across challenging
real-world scenarios including traffic violations, near-miss events, overtaking,
and U-turns, TUMTraf-V2X demonstrated that cooperative perception must operate
reliably across heterogeneous sensor configurations and complex driving manoeuvres
to be considered deployment-grade. Subsequent work has continued to push
deployment-grade evaluation further, with datasets such as
CATS-V2V~\cite{li2025catsv2v} targeting adverse weather and lighting conditions,
and UrbanIng-V2X~\cite{sekaran2025urbaningv2x} extending cooperative sensing
across multiple intersections with heterogeneous vehicle and infrastructure
agents.

\section{Open Challenges and Research Gaps}
\label{sec:gaps}

A critical evaluation of contemporary literature reveals unresolved challenges at
every level of the stereo cooperative perception stack. At the single-vehicle
level, stereo depth estimation remains unreliable in textureless regions and
degrades significantly with distance, while stereo-based 3D detection struggles
with severe occlusion and sparse geometric coverage. These upstream limitations
propagate directly into any downstream fusion pipeline, and are further compounded
by unresolved architectural gaps that restrict the transition of cooperative
perception from simulated environments to dependable real-world deployment. This
section synthesises the principal open challenges across the cooperative
perception stack.

\subsection{Communication Latency and Temporal Asynchrony}

The foundational assumption of multi-agent spatial coordination is that shared
perception data can be instantaneously mapped into a unified coordinate space. In
practical deployments, however, Vehicle-to-Vehicle (V2V) communication introduces
highly variable network propagation and queuing delays. Feature maps or finalised
bounding boxes transmitted with even a minor 200\,ms latency profile may
correspond to a physical spatial position that the transmitting vehicle has
already vacated. While frameworks such as SyncNet~\cite{lei2022} have attempted to
mitigate this temporal latency by estimating motion-compensated features to
predict object trajectories during network delays, the problem remains unresolved
under highly dynamic ego-motion. When vehicles undergo rapid manoeuvres such as
sudden braking, aggressive overtaking, or tight cornering at urban intersections,
idealised latency-compensation heuristics collapse, resulting in phantom
detections and duplicated tracks within the fused perception map.

\subsection{Pose Uncertainty and the Noisy-on-Noisy Alignment Problem}

Fusing spatial coordinates across a distributed network of moving vehicles
requires absolute precision regarding each agent's localised position within a
common global reference frame. Real-world Global Navigation Satellite System
(GNSS) localisation errors of even 1 to 2 metres can cause catastrophic geometric
misalignments in the aggregated tracking space. Although advanced frameworks like
CoAlign~\cite{lu2023} introduce robust pose-correction modules to optimise
geometric consistency, these systems suffer severe degradation in GNSS-degraded
environments such as deep urban canyons, underpasses, or tunnels.

This spatial alignment challenge is severely compounded within stereo-based
collaborative systems. Unlike active LiDAR sensors, which provide highly precise,
deterministic point-cloud measurements across a standard operating radius, stereo
vision systems carry inherent depth estimation noise that scales quadratically
with physical distance. Consequently, attempting to aggregate coordinate fields
from multiple stereo-equipped vehicles creates a compounded noisy-on-noisy spatial
alignment problem: because the baseline error models fluctuate independently
across each vehicle's unique perspective and distance to the target, standard
alignment techniques optimised for LiDAR-centric systems are structurally
unequipped to handle this geometric variance.

\subsection{Fleet Heterogeneity and Cross-Model Dependencies}

A restrictive bottleneck in contemporary cooperative perception is the systemic
reliance on software and hardware homogeneity across collaborating nodes. The vast
majority of state-of-the-art intermediate fusion networks require every
participating agent to execute identical deep learning model backbones, resolution
parameters, and feature extraction layers to ensure that the exchanged latent
tensors align semantically. In a realistic deployment scenario, autonomous vehicle
fleets are inherently heterogeneous, consisting of vehicles from different
manufacturers running entirely distinct perception setups. While decentralised
late fusion and head-level fusion models such as HEAD~\cite{qu2024} provide a
model-agnostic alternative, existing literature frequently trades architectural
flexibility for a reduction in localisation accuracy. Developing a cooperative
framework that maintains high spatial fidelity while seamlessly ingesting
detections from completely heterogeneous single-agent architectures remains an
open engineering challenge.

\subsection{Late Fusion Behaviour Under Stereo-Sourced Inputs}

Model-agnostic late fusion frameworks, including HEAD~\cite{qu2024} and prior
object-level aggregation methods, are evaluated exclusively on LiDAR-sourced
bounding boxes, whether drawn from real-world datasets such as
V2V4Real~\cite{xu2023v2v4real} and TUMTraf-V2X~\cite{zimmer2024}, or from
simulated environments such as OPV2V~\cite{xu2022opv2v}. LiDAR-sourced boxes carry
positional uncertainty that is largely range-independent within a sensor's
operating envelope, a property that clustering algorithms such as IoU matching and
Kalman filtering implicitly rely upon when resolving duplicate detections across
agents.

Stereo-sourced bounding boxes violate this assumption. Stereo depth error grows
quadratically with distance, meaning positional uncertainty in stereo-sourced
boxes is inherently range-dependent and perspective-specific across agents. When a
late fusion layer receives boxes from two stereo-equipped vehicles with different
baselines, focal lengths, and target distances, the noise profiles of the incoming
detections are structurally heterogeneous in a way that LiDAR-calibrated fusion
logic does not account for.

The empirical behaviour of late fusion pipelines operating strictly on
stereo-sourced inputs has not been characterised in the literature: how precision
degrades under range-dependent noise, how sensitive fusion is to inter-vehicle
baseline variation, and how detection quality evolves as target distance increases
all remain open questions. Establishing this baseline is a necessary step toward
designing cooperative perception systems that can operate effectively without
LiDAR infrastructure.

\bigskip
\noindent The three research domains surveyed in this chapter form a natural
end-to-end pipeline, yet their integration under a passive stereo sensing paradigm
leaves critical gaps at every stage: from the reliability of stereo depth under
real-world conditions, to the propagation of range-dependent noise through 3D
lifting, to the behaviour of cooperative fusion operating without LiDAR-sourced
inputs. The following chapter presents the system architecture developed to
address these gaps.
